\section{AIME I 2006}
        \begin{enumerate}
        	\item (\textit{AIME I 2006}) In quadrilateral $ABCD , \angle B$ is a right angle, diagonal $\overline{AC}$ is perpendicular to $\overline{CD},  AB=18, BC=21,$ and $CD=14.$ Find the perimeter of $ABCD.$


 
 \textbf{Answer: }084



\textbf{Solution 1}

We construct the following diagram: 
                    \begin{figure}[H]
                        \centering
                        \includegraphics[width=0.5\textwidth]{Solution/AIME_2006_I_1_Solution_Figure_1.png}
                    \end{figure}
                    Using the \href{https://artofproblemsolving.com/wiki/index.php?title=Pythagorean\_Theorem}{Pythagorean Theorem}, we get the following two equations:
 \[AD^2 = AC^2 + CD^2\]
 \[AC^2 = AB^2 + BC^2\]Substituting $AB^2 + BC^2$ for $AC^2$ gives us $AD^2 = AB^2 + BC^2 + CD^2$. Plugging in the given information, we get $AD^2 = 18^2 + 21^2 + 14^2 = 961 \implies AD = 31$, so the perimeter is $AB+BC+CD+AD = 18+21+14+31 = \boxed{084}$.

 




	\item (\textit{AIME I 2006}) Let set $\mathcal{A}$ be a 90-element subset of $\{1,2,3,\ldots,100\},$ and let $S$ be the sum of the elements of $\mathcal{A}.$ Find the number of possible values of $S.$


 
 \textbf{Answer: }901



\textbf{Solution}

The smallest $S$ is $1+2+ \ldots +90 = 91 \cdot 45 = 4095$. The largest $S$ is $11+12+ \ldots +100=111\cdot 45=4995$. All numbers between $4095$ and $4995$ are possible values of S, so the number of possible values of S is $4995-4095+1=\boxed{901}$.

 Alternatively, for ease of calculation, let set $\mathcal{B}$ be a 10-element subset of $\{1,2,3,\ldots,100\}$, and let $T$ be the sum of the elements of $\mathcal{B}$. Note that the number of possible $S$ is the number of possible $T=5050-S$. The smallest possible $T$ is $1+2+ \ldots +10 = 55$ and the largest is $91+92+ \ldots + 100 = 955$, so the number of possible values of T, and therefore S, is $955-55+1=\boxed{901}$.

 




	\item (\textit{AIME I 2006}) Find the least positive integer such that when its leftmost digit is deleted, the resulting integer is $1/29$ of the original integer.


 
 \textbf{Answer: }725



\textbf{Solution 1}

Suppose the original number is $N = \overline{a_na_{n-1}\ldots a_1a_0},$ where the $a_i$ are digits and the first digit, $a_n,$ is nonzero. Then the number we create is $N_0 = \overline{a_{n-1}\ldots a_1a_0},$ so 
 \[N = 29N_0.\] But $N$ is $N_0$ with the digit $a_n$ added to the left, so $N = N_0 + a_n \cdot 10^n.$ Thus, 
 \[N_0 + a_n\cdot 10^n = 29N_0\] 
 \[a_n \cdot 10^n = 28N_0.\] The right-hand side of this equation is divisible by seven, so the left-hand side must also be divisible by seven. The number $10^n$ is never divisible by $7,$ so $a_n$ must be divisible by $7.$ But $a_n$ is a nonzero digit, so the only possibility is $a_n = 7.$ This gives 
 \[7 \cdot 10^n = 28N_0\] or 
 \[10^n = 4N_0.\] Now, we want to minimize \textit{both} $n$ and $N_0,$ so we take $N_0 = 25$ and $n = 2.$ Then 
 \[N = 7 \cdot 10^2 + 25 = \boxed{725},\] and indeed, $725 = 29 \cdot 25.$ $\square$

 





\textbf{Solution 2}

Let $N$ be the required number, and $N'$ be $N$ with the first digit deleted. Now, we know that $N<1000$ (because this is an AIME problem). Thus, $N$ has $1,$ $2$ or $3$ digits. Checking the other cases, we see that it must have $3$ digits.Let $N=\overline{abc}$, so $N=100a+10b+c$. Thus, $N'=\overline{bc}=10b+c$. By the constraints of the problem, we see that $N=29N'$, so 
 \[100a+10b+c=29(10b+c).\] Now, we subtract and divide to get 
 \[100a=28(10b+c)\] 
 \[25a=70b+7c.\] Clearly, $c$ must be a multiple of $5$ because both $25a$ and $70b$ are multiples of $5$. Thus, $c=5$. Now, we plug that into the equation: 
 \[25a=70b+7(5)\] 
 \[25a=70b+35\] 
 \[5a=14b+7.\] By the same line of reasoning as earlier, $a=7$. We again plug that into the equation to get 
 \[35=14b+7\] 
 \[b=2.\] Now, since $a=7$, $b=2$, and $c=5$, our number $N=100a+10b+c=\boxed{725}$.

 Here's another way to finish using this solution. From the above, you have 
 \[100a = 28(10b + c).\] Divide by $4$, and you get 
 \[25a = 7(10b + c).\] This means that $25a$ has to be divisible by $7$, and hence $a = 7.$ Now, solve for $25 = 10b + c$, which gives you $a = 7, b = 2, c = 5$, giving you the number $\boxed{725}$

 





\textbf{Solution 3 (Very fast)}

Assume the answer is a three digit number (because it has to be 3 or less and 2 is way too small). Let b be the last two digits (and so less than 100), and let a be the first digit (and so a multiple of 100). Then we have \( a+b=29b \implies a=28b \implies 4 | a, 7 | a \). Since a is a multiple of 100 it is automatically divisible by 4, and so we pick \(a=700\) to satisfy the other condition, yielding \(b=25\) and so our answer is \( \boxed{725} \).~michael

 





\textbf{Solution 4 (Quick)}

Note that if we let the last digit be $c$ we must have $9c \equiv c \pmod{10}.$ Thus we either have $c=0$ which we can quickly check to be impossible (since the number after digit removal could be 10,20,30) or $c=5.$ Testing 5, 15, and 25 as the numbers after removal we find that our answer is clearly $29 \cdot 25 = 725.$

 (Note that the quick checking of six numbers was possible thanks to AIME problems having answers less than 1000).

 





\textbf{Solution 5}

First we try $2$ digit numbers like $\overline{ab}$. Removing the leftmost digit gives us $b$. Now, using the information we are given we can write the equation: 
 \[\dfrac{b}{10a+b} = \dfrac{1}{29} \Longrightarrow 10a+b=29b \Longrightarrow 5a = 14b\]. Since $5$ and $14$ are relatively prime, the smallest case that works is $a = 14$, and $b = 5$, but $a$ is an digit so the number must be at least $3$ digits. We let the number be $\overline{abc}$. So, using the info we know, we can conclude that 
 \[\dfrac{10b+c}{100a+10b+c} = \dfrac{1}{29} \Longrightarrow 25a-70b-7c = 0\]. Rerraranging gives us 
 \[25a = 7(10b+c)\]. So, sicne $25$ and $7$ are relatively prime, the lowest possible case is that $a = 7$, and $10b+c = 25$.So, we have $b = 2, c = 5$. So the desired number is $\boxed{725}$

 -jb2015007

 




	\item (\textit{AIME I 2006}) Let $N$ be the number of consecutive 0's at the right end of the decimal representation of the product $1!2!3!4!\cdots99!100!.$ Find the remainder when $N$ is divided by 1000.


 
 \textbf{Answer: }124



\textbf{Solution}

A number in decimal notation ends in a zero for each power of ten which divides it.  Thus, we need to count both the number of 5s and the number of 2s dividing into our given expression.  Since there are clearly more 2s than 5s, it is sufficient to count the number of 5s.

 One way to do this is as follows: $96$ of the numbers $1!,\ 2!,\ 3!,\ 100!$ have a factor of $5$. $91$ have a factor of $10$. $86$ have a factor of $15$. And so on. This gives us an initial count of $96 + 91 + 86 + \ldots + 1$.  Summing this \href{https://artofproblemsolving.com/wiki/index.php?title=Arithmetic\_series}{arithmetic series} of $20$ terms, we get $970$.  However, we have neglected some powers of $5$ - every $n!$ term for $n\geq25$ has an additional power of $5$ dividing it, for $76$ extra; every n! for $n\geq 50$ has one more in addition to that, for a total of $51$ extra; and similarly there are $26$ extra from those larger than $75$ and $1$ extra from $100$.  Thus, our final total is $970 + 76 + 51 + 26 + 1 = 1124$, and the answer is $\boxed{124}$.

 





\textbf{Solution 2 (Legendre's Formula )}

This problem can be easily solved using \href{https://artofproblemsolving.com/wiki/index.php?title=Legendre\%27s\_Formula}{Legendre's Formula}:

 First is to account for all of the factorials that are greater than $5n$ or $5, 10, 15, 20...100$, then the factorials that are greater than $5^2n$ or $25, 50, 75, 100$. Since $5^3=125>100$ we do not need to account for it.

 This gives $96+91+86+...+1$ and $76+51+26+1$ respectively. Adding all of these terms up gives 1124 or $\boxed{124}$.

 




	\item (\textit{AIME I 2006}) The number $\sqrt{104\sqrt{6}+468\sqrt{10}+144\sqrt{15}+2006}$ can be written as $a\sqrt{2}+b\sqrt{3}+c\sqrt{5},$ where $a, b,$ and $c$ are positive integers. Find $abc.$


 
 \textbf{Answer: }936



\textbf{Solution 1}

We begin by \href{https://artofproblemsolving.com/wiki/index.php?title=Equate}{ equating} the two expressions:

 
 \[a\sqrt{2}+b\sqrt{3}+c\sqrt{5} = \sqrt{104\sqrt{6}+468\sqrt{10}+144\sqrt{15}+2006}\]

 Squaring both sides yields: 

 
 \[2ab\sqrt{6} + 2ac\sqrt{10} + 2bc\sqrt{15} + 2a^2 + 3b^2 + 5c^2 = 104\sqrt{6}+468\sqrt{10}+144\sqrt{15}+2006\] 

 Since $a$, $b$, and $c$ are integers, we can match coefficients: 

 
 \begin{align*} 2ab\sqrt{6} &= 104\sqrt{6} \\  2ac\sqrt{10} &=468\sqrt{10} \\  2bc\sqrt{15} &=144\sqrt{15}\\  2a^2 + 3b^2 + 5c^2 &=2006 \end{align*}
 

 Solving the first three equations gives: 
 \begin{eqnarray*}ab &=& 52\\  ac &=& 234\\  bc &=& 72 \end{eqnarray*} 

 Multiplying these equations gives $(abc)^2 = 52 \cdot 234 \cdot 72 = 2^63^413^2 \Longrightarrow abc = \boxed{936}$. 

 


 





\textbf{Solution 2}

We realize that the quantity under the largest radical is a perfect square and attempt to rewrite the radicand as a square. Start by setting $x=\sqrt{2}$, $y=\sqrt{3}$, and $z=\sqrt{5}$.  Since

 
 \[(px+qy+rz)^2=p^2x^2+q^2y^2+r^2z^2+2(pqxy+prxz+qryz)\]

 we attempt to rewrite the radicand in this form:

 
 \[2006+2(52xy+234xz+72yz)\]

 Factoring, we see that $52=13\cdot4$, $234=13\cdot18$, and $72=4\cdot18$. Setting $p=13$, $q=4$, and $r=18$, we see that

 
 \[2006=13^2x^2+4^2y^2+18^2z^2=169\cdot2+16\cdot3+324\cdot5\]

 so our numbers check. Thus $104\sqrt{2}+468\sqrt{3}+144\sqrt{5}+2006=(13\sqrt{2}+4\sqrt{3}+18\sqrt{5})^2$. Square rooting gives us $13\sqrt{2}+4\sqrt{3}+18\sqrt{5}$ and our answer is $13\cdot4\cdot18=\boxed{936}$

 




	\item (\textit{AIME I 2006}) Let $\mathcal{S}$ be the set of real numbers that can be represented as repeating decimals of the form $0.\overline{abc}$ where $a, b, c$ are distinct digits. Find the sum of the elements of $\mathcal{S}.$


 
 \textbf{Answer: }360



\textbf{Solution 1}

Numbers of the form $0.\overline{abc}$ can be written as $\frac{abc}{999}$. There are $10\times9\times8=720$ such numbers. Each digit will appear in each place value $\frac{720}{10}=72$ times, and the sum of the digits, 0 through 9, is 45. So the sum of all the numbers is $\frac{45\times72\times111}{999}= \boxed{360}$.

 





\textbf{Solution 2}

Alternatively, for every number, $0.\overline{abc}$, there will be exactly one other number, such that when they are added together, the sum is $0.\overline{999}$, or, more precisely, 1. As an example, $.\overline{123}+.\overline{876}=.\overline{999} \Longrightarrow 1$.

 Thus, the solution can be determined by dividing the total number of \href{https://artofproblemsolving.com/wiki/index.php?title=Permutation}{permutations} by 2. The answer is $\frac{10 \cdot 9 \cdot 8}{2} = \frac{720}{2}= \boxed{360}$.

 Another method, albeit a little risky, that can be used is to note that the numbers between 1 and 999 with distinct digits average out to $\frac{999}{2}$. Then the total sum becomes $\frac{\frac{999}{2}\times720}{999}$ which reduces to $\boxed{360}$

 





\textbf{Solution 3}

Another way is to do complementary counting and casework. $0.\overline{abc} = abc/999$. Excluding numbers with repeating digits our answer would just be $\frac{999 \cdot (1000/2)}{999} = 500$. We subtract off the sum of the numbers with at least two digits being the same:

 
Case 1: $aba$: Sum = $1010(1 + 2 + 3 + \dots + 9) + 450\cdot 10 = 999 \cdot 50$.

 
Case 2: $aab$: Sum = $1100(1 + 2 + 3 + \dots + 9) + 45\cdot 10 = 999 \cdot 50$

 
Case 3: $baa$: Sum = $110(1 + 2 + 3 + \dots + 9) + 4500\cdot 10 = 999 \cdot 50$.

 
We overcounted the case where $a = b$ \textbf{twice} (the numbers with all 3 digits being the same). The sum of these numbers is $111(1 + 2 + 3 + \dots + 9) = 999 \cdot 5$. 

 
So, our final answer is $\frac{999 \cdot 500 - 999 \cdot 150 + 999 \cdot 5 \cdot 2}{999} = \boxed{360}$.

 




	\item (\textit{AIME I 2006}) An angle is drawn on a set of equally spaced parallel lines as shown. The ratio of the area of shaded region $C$ to the area of shaded region $B$ is $\frac{11}{5}$. Find the ratio of shaded region $D$ to the area of shaded region $A$.


 
                    \begin{figure}[H]
                        \centering
                        \includegraphics[width=0.5\textwidth]{Problem/AIME_2006_I_7_Problem_Figure_1.png}
                    \end{figure}
                    


 
 \textbf{Answer: }408



\textbf{Solution 1}

Note that the apex of the angle is not on the parallel lines. Set up a \href{https://artofproblemsolving.com/wiki/index.php?title=Coordinate\_proof&action=edit&redlink=1}{coordinate proof}.

 Let the set of parallel lines be \href{https://artofproblemsolving.com/wiki/index.php?title=Perpendicular}{perpendicular} to the \href{https://artofproblemsolving.com/wiki/index.php?title=X-axis}{x-axis}, such that they cross it at $0, 1, 2 \ldots$. The base of region $\mathcal{A}$ is on the line $x = 1$. The bigger base of region $\mathcal{D}$ is on the line $x = 7$. Let the top side of the angle be $y = x - s$ and the bottom side be x-axis, as dividing the angle doesn't change the problem.

 Since the area of the triangle is equal to $\frac{1}{2}bh$,

 
 \[\frac{\textrm{Region\ }\mathcal{C}}{\textrm{Region\ }\mathcal{B}} = \frac{11}{5} = \frac{\frac 12(5-s)^2 - \frac 12(4-s)^2}{\frac 12(3-s)^2 - \frac12(2-s)^2}\]

 Solve this to find that $s = \frac{5}{6}$.

 Using the same reasoning as above, we get $\frac{\textrm{Region\ }\mathcal{D}}{\textrm{Region\ }\mathcal{A}} = \frac{\frac 12(7-s)^2 - \frac 12(6-s)^2}{\frac 12(1-s)^2}$, which is $\boxed{408}$.

 





\textbf{Solution 2}

Note that the sections between the two transversals can be divided into one small triangle and a number of trapezoids. Let one side length (not on a parallel line) of the small triangle be $x$ and the area of it be $x^2$. Also, let all sections of the line on the same side as the side with length $x$ on a trapezoid be equal to $1$. 

 Move on to the second-smallest triangle, formed by attaching this triangle with the next trapezoid. Parallel lines give us similar triangles, so we know the proportion of this triangle to the previous triangle is ${\left(\frac{x+1}{x}\right)}^2$. Multiplying, we get $(x+1)^2$ as the area of the triangle, so the area of the trapezoid is $2x+1$. Repeating this process, we get that the area of B is $2x+3$, the area of C is $2x+7$, and the area of D is $2x+11$.

 We can now use the given condition that the ratio of C and B is $\frac{11}{5}$.

 $\frac{11}{5} = \frac{2x+7}{2x+3}$ gives us $x = \frac{1}{6}$

 So now we compute the ratio of D and A, which is $\frac{(2)(\frac{1}{6}) + 11}{(\frac{1}{6})^2} = \boxed{408.}$

 Edit: fixed misplaced brackets

 





\textbf{Solution 3 (Bash)}

Let the distances from the apex to the parallel lines be $x$ and $y$ and the distance between the intersections be $a,b.$ We know the area ratio means $\frac{(x+4a)(y+4b)-(x+3a)(y+3b)}{(x+2a)(y+2b)-(x+a)(y+b)} =\frac{5}{11}$ which simplifying yields $ab = 3ay+3bx.$ The ratio we seek is $\frac{(x+6a)(y+6b)-(x+5a)(y+5b)}{xy} =\frac{ay+yx+11ab}{xy}.$ We know that $ab = 3ay+3bx$ so the ratio we seed is $\frac{33(ay+yx)}{11xy}.$ Finally note that by similar triangles $\frac{x}{x+a} =\frac{y}{y+b} \implies bx = ya.$ Therefore the ratio we seek is $\frac{66(ay)}{11xy} =\frac{66a}{11x}.$ Finally note that $ab=3ay+3bx \implies ab = 6bx \implies a = 6x$ so the final ratio is $6 \cdot 68 = \boxed{408}.$

 




	\item (\textit{AIME I 2006}) Hexagon $ABCDEF$ is divided into five rhombuses, $P, Q, R, S,$ and $T$, as shown. Rhombuses $P, Q, R,$ and $S$ are congruent, and each has area $\sqrt{2006}.$ Let $K$ be the area of rhombus $T$. Given that $K$ is a positive integer, find the number of possible values for $K.$


 
                    \begin{figure}[H]
                        \centering
                        \includegraphics[width=0.5\textwidth]{Problem/AIME_2006_I_8_Problem_Figure_1.png}
                    \end{figure}
                    


 
 \textbf{Answer: }089



\textbf{Solution 1}

Let $x$ denote the common side length of the rhombi, and let $y$ denote one of the smaller interior \href{https://artofproblemsolving.com/wiki/index.php?title=Angle}{angles} of rhombus $\mathcal{P}$. Then $x^2\sin(y)=\sqrt{2006}$.  We also see that $K=x^2\sin(2y) \implies K=2x^2\sin y \cdot \cos y \implies K = 2\sqrt{2006}\cdot \cos y$.  Thus $K$ can be any positive integer in the \href{https://artofproblemsolving.com/wiki/index.php?title=Interval}{interval} $(0, 2\sqrt{2006})$. Since $2\sqrt{2006} = \sqrt{8024}$ and $89^2 = 7921 < 8024 < 8100 = 90^2$, $K$ can be any \href{https://artofproblemsolving.com/wiki/index.php?title=Integer}{integer} between $1$ and $89$, inclusive, so the number of positive values for $K$ is $\boxed{089}$.

 


 





\textbf{Solution 2}

Call the side length of each rhombus $w$. $w$ is the width of the rhombus. Call the height h, where $wh=\sqrt{2006}$. The height of rhombus T would be 2h, and the width would be $\sqrt{w^2-h^2}$. Substitute the first equation to get $\sqrt{\frac{2006}{h^2}-h^2}$. Then the area of the rhombus would be $2h * \sqrt{\frac{2006}{h^2}-h^2}$. Combine like terms to get $2 * \sqrt{2006-h^4}$. This expression equals an integer K. $2006-h^4$ specifically must be in the form $n^2/4$. There is no restriction on h as long as it is a positive real number, so all we have to do is find all the positive possible values of $n^2$ for $2006-h^4$. Now, quick testing shows that $44^2 < 2006$ and $45^2>2006$, but we must also test $44.5^2$, because the product of two will make it an integer. $44.5^2$ is also less than $2006$, so we have numbers 1-44, times two because 0.5 can be added to each of time, plus 1, because 0.5 is also a valid value. (notice 0 is not valid because the height must be a positive number) That gives us $44*2+1=$ $\boxed{089}$

 -jackshi2006

 





\textbf{Solution 3}


                    \begin{figure}[H]
                        \centering
                        \includegraphics[width=0.5\textwidth]{Solution/AIME_2006_I_8_Solution_Figure_1.png}
                    \end{figure}
                    To determine the possible values of $[GIHJ],$ we must determine the maximum and minimum possible areas.

 In the case where the $4$ rhombi are squares, we have $[GIHJ]=0,$ implying the minimum possible positive-integer-valued area is $1.$

 Denote the length $HC=a$ and $KH=b.$ We have
 \[KI=\sqrt{a^2-b^2}\]by the Pythagorean Theorem, which implies
 \[[IBCH]=a\sqrt{a^2-b^2}=\sqrt{2006}\]and
 \[[GIHJ]=2b\sqrt{a^2-b^2}.\]The first equation yields
 \[\sqrt{a^2-b^2}=\frac{\sqrt{2006}}{a}.\]Plugging into the second, we have
 \[[GIHJ]=2\sqrt{2006}\frac{b}{a}.\]The maximal value of $\frac{b}{a}$ occurs when the height of $ABCDEF$ is minimized, which means
 \[\frac{b}{a}\leq 1.\]Plugging back up, we have
 \[[GIHJ]\leq 2\sqrt{2006}=\sqrt{8024}.\]We have
 \[\lfloor \sqrt{8024} \rfloor = 89,\]thus our answer is
 \[89-1+1=\boxed{089}.\]

 




	\item (\textit{AIME I 2006}) The sequence $a_1, a_2, \ldots$ is geometric with $a_1=a$ and common ratio $r,$ where $a$ and $r$ are positive integers. Given that $\log_8 a_1+\log_8 a_2+\cdots+\log_8 a_{12} = 2006,$ find the number of possible ordered pairs $(a,r).$


 
 \textbf{Answer: }046



\textbf{Solution 1}


 \[\log_8 a_1+\log_8 a_2+\ldots+\log_8 a_{12}= \log_8 a+\log_8 (ar)+\ldots+\log_8 (ar^{11}) \\ = \log_8(a\cdot ar\cdot ar^2\cdot \cdots \cdot ar^{11}) = \log_8  (a^{12}r^{66})\]

 So our question is equivalent to solving $\log_8 (a^{12}r^{66})=2006$ for $a, r$ \href{https://artofproblemsolving.com/wiki/index.php?title=Positive\_integer}{positive integers}. $a^{12}r^{66}=8^{2006} = (2^3)^{2006} = (2^6)^{1003}$ so $a^{2}r^{11}=2^{1003}$.

 The product of $a^2$ and $r^{11}$ is a power of 2.  Since both numbers have to be integers, this means that $a$ and $r$ are themselves powers of 2.  Now, let $a=2^x$ and $r=2^y$:

 
 \begin{eqnarray*}(2^x)^2\cdot(2^y)^{11}&=&2^{1003}\\ 2^{2x}\cdot 2^{11y}&=&2^{1003}\\ 2x+11y&=&1003\\ y&=&\frac{1003-2x}{11} \end{eqnarray*}

 For $y$ to be an integer, the 

 




	\item (\textit{AIME I 2006}) Eight circles of diameter 1 are packed in the first quadrant of the coordinate plane as shown. Let region $\mathcal{R}$ be the union of the eight circular regions. Line $l,$ with slope 3, divides $\mathcal{R}$ into two regions of equal area. Line $l$'s equation can be expressed in the form $ax=by+c,$ where $a, b,$ and $c$ are positive integers whose greatest common divisor is 1. Find $a^2+b^2+c^2.$


 
                    \begin{figure}[H]
                        \centering
                        \includegraphics[width=0.5\textwidth]{Problem/AIME_2006_I_10_Problem_Figure_1.png}
                    \end{figure}
                    


 
 \textbf{Answer: }065



\textbf{Solutions}


                    \begin{figure}[H]
                        \centering
                        \includegraphics[width=0.5\textwidth]{Solution/AIME_2006_I_10_Solution_Figure_1.png}
                    \end{figure}
                    

 





\textbf{Solution 1}

The line passing through the \href{https://artofproblemsolving.com/wiki/index.php?title=Tangent\_(geometry)}{tangency point} of the bottom left circle and the one to its right and through the tangency of the top circle in the middle column and the one beneath it is the line we are looking for: a line passing through the tangency of two circles cuts congruent areas, so our line cuts through the four aforementioned circles splitting into congruent areas, and there are an additional two circles on each side. The line passes through $\left(1,\frac 12\right)$ and $\left(\frac 32,2\right)$, which can be easily solved to be $6x = 2y + 5$. Thus, $a^2 + b^2 + c^2 = \boxed{065}$.

 





\textbf{Solution 2}

Assume that if unit \href{https://artofproblemsolving.com/wiki/index.php?title=Square}{squares} are drawn circumscribing the circles, then the line will divide the area of the \href{https://artofproblemsolving.com/wiki/index.php?title=Concave}{concave} hexagonal region of the squares equally (as of yet, there is no substantiation that such would work, and definitely will not work in general). Denote the intersection of the line and the \href{https://artofproblemsolving.com/wiki/index.php?title=X-axis}{x-axis} as $(x, 0)$. 

 The line divides the region into 2 sections. The left piece is a \href{https://artofproblemsolving.com/wiki/index.php?title=Trapezoid}{trapezoid}, with its area $\frac{1}{2}((x) + (x+1))(3) = 3x + \frac{3}{2}$. The right piece is the addition of a \href{https://artofproblemsolving.com/wiki/index.php?title=Trapezoid}{trapezoid} and a \href{https://artofproblemsolving.com/wiki/index.php?title=Rectangle}{rectangle}, and the areas are $\frac{1}{2}((1-x) + (2-x))(3)$ and $2 \cdot 1 = 2$, totaling $\frac{13}{2} - 3x$. Since we want the two regions to be equal, we find that $3x + \frac 32 = \frac {13}2 - 3x$, so $x = \frac{5}{6}$. 

 We have that $\left(\frac 56, 0\right)$ is a point on the line of slope 3, so $y - 0 = 3\left(x - \frac 56\right) \Longrightarrow 6x = 2y + 5$. Our answer is $2^2 + 5^2 + 6^2 = 65$.

 We now assess the validity of our starting assumption. We can do that by seeing that our answer passes through the tangency of the two circles, cutting congruent areas, a result explored in solution 1.

 


 





\textbf{Solution 3}

This problem looks daunting at a first glance, but we can make geometric inequality inferences by drawing lines that simplify the problem by removing sections of the total area. To begin, we can eliminate the possibility of the line intersecting the circle on the top left (call it circle A), or the circle on the bottom right (call it circle B). This is can be seen visually by drawing a line with slope 3 that is tangent to either of these circles. The area is clearly larger on one side; this can be proven by counting full circles. We can go on with the same mindset and eliminate the circle below circle A and the circle above circle B. By removing pairs of circles and proving the line will never intersect with them, we can safely work with whatever is remaining. By now you should have 4 circles making an L shape (waluigi style). Now the two biggest contenders for this method are the two circles on the bottom row. Using the same strategy, we can see that a line that goes through the tangent point of these two circles also goes through the tangent point of the other two circles. This clearly will cut the 4 circles into two regions of equal area. Using linear algebra we get: $6x = 2y + 5$. The answer is then $6^2 + 2^2 + 5^2 = \boxed{065}$. This solution is an alternate explanation to solution 1.

 
-jackshi2006

 




	\item (\textit{AIME I 2006}) A collection of 8 cubes consists of one cube with edge-length $k$ for each integer $k, 1 \le k \le 8.$ A tower is to be built using all 8 cubes according to the rules: 


 
                \begin{itemize}
                    \item Any cube may be the bottom cube in the tower.
\item The cube immediately on top of a cube with edge-length $k$ must have edge-length at most $k+2.$
                \end{itemize}


                Let $T$ be the number of different towers than can be constructed. What is the remainder when $T$ is divided by 1000?


 
 \textbf{Answer: }458



\textbf{Solution}

We proceed \href{https://artofproblemsolving.com/wiki/index.php?title=Recursion}{ recursively}.  Suppose we can build $T_m$ towers using blocks of size $1, 2, \ldots, m$.  How many towers can we build using blocks of size $1, 2, \ldots, m, m + 1$?  If we remove the block of size $m + 1$ from such a tower (keeping all other blocks in order), we get a valid tower using blocks $1, 2, \ldots, m$.  Given a tower using blocks $1, 2, \ldots, m$ (with $m \geq 2$), we can insert the block of size $m + 1$ in exactly 3 places: at the beginning, immediately following the block of size $m - 1$ or immediately following the block of size $m$.  Thus, there are 3 times as many towers using blocks of size $1, 2, \ldots, m, m + 1$ as there are towers using only $1, 2, \ldots, m$.  There are 2 towers which use blocks $1, 2$, so there are $2\cdot 3^6 = 1458$ towers using blocks $1, 2, \ldots, 8$, so the answer is $\boxed{458}$.

 (Note that we cannot say, "there is one tower using the block $1$, so there are $3^7$ towers using the blocks $1, 2, \ldots, 8$."  The reason this fails is that our recursion only worked when $m \geq 2$: when $m = 1$, there are only 2 places to insert a block of size $m + 1 = 2$, at the beginning or at the end, rather than the 3 places we have at later stages.  Also, note that this method generalizes directly to seeking the number of towers where we change the second rule to read, "The cube immediately on top of a cube with edge-length $k$ must have edge-length at most $k + n$," where $n$ can be any fixed integer.)

 




	\item (\textit{AIME I 2006}) Find the sum of the values of $x$ such that $\cos^3 3x+ \cos^3 5x = 8 \cos^3 4x \cos^3 x,$ where $x$ is measured in degrees and $100< x< 200.$


 
 \textbf{Answer: }906



\textbf{Solution 1 (Pure trig)}

Observe that $2\cos 4x\cos x = \cos 5x + \cos 3x$ by the sum-to-product formulas. Defining $a = \cos 3x$ and $b = \cos 5x$, we have $a^3 + b^3 = (a+b)^3 \rightarrow ab(a+b) = 0$. But $a+b = 2\cos 4x\cos x$, so we require $\cos x = 0$, $\cos 3x = 0$, $\cos 4x = 0$, or $\cos 5x = 0$.

 Hence we see by careful analysis of the cases that the solution set is $A = \{150, 126, 162, 198, 112.5, 157.5\}$ and thus $\sum_{x \in A} x = \boxed{906}$.

 





\textbf{Solution 2 (Complex trig)}

Let $z=\cos(x)+i \sin(x)$. Then, we have that $|z|=1$, so $\overline{z}=\frac{1}{z}$. Thus, $\frac{1}{z}=\cos(x)-i \sin(x)$, and we have $\cos(x)=\frac{z+\frac{1}{z}}{2}$. Now, this means we also have

 
 \[z^{3}=\cos(3x)+i\sin(3x) \rightarrow z^{3}+\frac{1}{z^{3}} = 2\cos(3x) \rightarrow \frac{z^{3}+\frac{1}{z^{3}}}{2}=\cos(3x)\]

 
 \[\frac{z^{5}+\frac{1}{z^{5}}}{2}=\cos(5x)\]

 Then, we seek the solutions to 

 
 \[\frac{1}{8} \cdot \left(z^{3}+\frac{1}{z^{3}}\right)^{3} + \frac{1}{8} \cdot \left(z^{5}+\frac{1}{z^{5}}\right)^{3} = 8 \cdot \frac{1}{64} \cdot \left(z+\frac{1}{z}\right)^{3}\left(z^{4}+\frac{1}{z^{4}}\right)^{3}\]

 
 \[\rightarrow \left(z^{3}+\frac{1}{z^{3}}\right)^{3}+\left(z^{5}+\frac{1}{z^{5}}\right)^{3}=\left(z^{3}+\frac{1}{z^{3}}+z^{5}+\frac{1}{z^{5}}\right)^{3}\]

 Now, if we let $a=z^{3}+\frac{1}{z^{3}}$ and we let $b=z^{5}+\frac{1}{z^{5}}$, we actually obtain the equation

 
 \[a^{3}+b^{3}=(a+b)^{3} \rightarrow ab(a+b)=0\]

 which means we either have

 
 \[\textbf{Case 1:} \hspace{3pt} z^{3}+\frac{1}{z^{3}}=0\]
 \[\textbf{Case 2:} \hspace{3pt} z^{5}+\frac{1}{z^{5}}=0\]
 \[\textbf{Case 3:} \hspace{5pt} z^{3}+\frac{1}{z^{3}}=z^{5}=\frac{1}{z^{5}}\]

 In Case 1, we have $z^{12}=1$ but $z^{6} \neq 1$. Thus, we have that the possibilities for $\arg[z]$ are $30^\circ$, $90^\circ$, $\underline{150}^\circ$, $210^\circ$, $270^\circ$, and $330^\circ$. 

 In Case 2, we have $z^{20}=1$ but $z^{10} \neq 1$. Thus, we have that the set of possible $z$ are 10th roots of unity shifted by $18^\circ$ counterclockwise, so we have

 
 \[\arg[z] \in \{18^\circ, 54^\circ, 90^\circ, \underline{126}^\circ, \underline{162}^\circ, \underline{198}^\circ, \dots, 342^\circ \}\]

 so we obtain $3$ values. In case 3, we have

 
 \[z^{3}+\frac{1}{z^{3}}=z^{5}+\frac{1}{z^{5}} \rightarrow z^{10}+z^{8}+z^{2}+1=0 \rightarrow (z^{8}+1)(z^{2}+1)=0\]

 so either $z \in \{i, -i\}$ or $z$ is a primitive 8th root of unity shifted by $22.5^\circ$, which gives 

 
 \[\arg[z] \in \{22.5^\circ, 67.5^\circ, \underline{112.5}^\circ, \underline{157.5}^\circ, 202.5^\circ, 247.5^\circ, 292.5^\circ, 337.5^\circ \}\]

 The sum of all underlined values is $150+126+162+198+112.5+157.5=\fbox{906}$

 




	\item (\textit{AIME I 2006}) For each even positive integer $x,$ let $g(x)$ denote the greatest power of 2 that divides $x.$ For example, $g(20)=4$ and $g(16)=16.$ For each positive integer $n,$ let $S_n=\sum_{k=1}^{2^{n-1}}g(2k).$ Find the greatest integer $n$ less than 1000 such that $S_n$ is a perfect square.


 
 \textbf{Answer: }899



\textbf{Solution 1}

Given $g : x \mapsto \max_{j : 2^j | x} 2^j$, consider $S_n = g(2) + \cdots + g(2^n)$. Define $S = \{2, 4, \ldots, 2^n\}$. There are $2^0$ elements of $S$ that are divisible by $2^n$, $2^1 - 2^0 = 2^0$ elements of $S$ that are divisible by $2^{n-1}$ but not by $2^n, \ldots,$ and $2^{n-1}-2^{n-2} = 2^{n-2}$ elements of $S$ that are divisible by $2^1$ but not by $2^2$.

 Thus 
 \begin{align*} S_n &= 2^0\cdot2^n + 2^0\cdot2^{n-1} + 2^1\cdot2^{n-2} + \cdots + 2^{n-2}\cdot2^1\\ &= 2^n + (n-1)2^{n-1}\\ &= 2^{n-1}(n+1).\end{align*}
 Let $2^k$ be the highest power of $2$ that divides $n+1$. Thus by the above formula, the highest power of $2$ that divides $S_n$ is $2^{k+n-1}$. For $S_n$ to be a perfect square, $k+n-1$ must be even. If $k$ is odd, then $n+1$ is even, hence $k+n-1$ is odd, and $S_n$ cannot be a perfect square. Hence $k$ must be even. In particular, as $n<1000$, we have five choices for $k$, namely $k=0,2,4,6,8$.

 If $k=0$, then $n+1$ is odd, so $k+n-1$ is odd, hence the largest power of $2$ dividing $S_n$ has an odd exponent, so $S_n$ is not a perfect square.

 In the other cases, note that $k+n-1$ is even, so the highest power of $2$ dividing $S_n$ will be a perfect square. In particular, $S_n$ will be a perfect square if and only if $(n+1)/2^{k}$ is an odd perfect square.

 If $k=2$, then $n<1000$ implies that $\frac{n+1}{4} \le 250$, so we have $n+1 = 4, 4 \cdot 3^2, \ldots, 4 \cdot 13^2, 4\cdot 3^2 \cdot 5^2$.

 If $k=4$, then $n<1000$ implies that $\frac{n+1}{16} \le 62$, so $n+1 = 16, 16 \cdot 3^2, 16 \cdot 5^2, 16 \cdot 7^2$.

 If $k=6$, then $n<1000$ implies that $\frac{n+1}{64}\le 15$, so $n+1=64,64\cdot 3^2$.

 If $k=8$, then $n<1000$ implies that $\frac{n+1}{256}\le 3$, so $n+1=256$.

 Comparing the largest term in each case, we find that the maximum possible $n$ such that $S_n$ is a perfect square is $4\cdot 3^2 \cdot 5^2 - 1 = \boxed{899}$.

 





\textbf{Solution 2}

First note that $g(k)=1$ if $k$ is odd and $2g(k/2)$ if $k$ is even.so $S_n=\sum_{k=1}^{2^{n-1}}g(2k). = \sum_{k=1}^{2^{n-1}}2g(k) = 2\sum_{k=1}^{2^{n-1}}g(k) = 2\sum_{k=1}^{2^{n-2}} g(2k-1)+g(2k).$$2k-1$ must be odd so this reduces to $2\sum_{k=1}^{2^{n-2}}1+g(2k) = 2(2^{n-2}+\sum_{k=1}^{2^{n-2}}g(2k)).$  Thus $S_n=2(2^{n-2}+S_{n-1})=2^{n-1}+2S_{n-1}.$ Further noting that $S_0=1$ we can see that $S_n=2^{n-1}\cdot (n-1)+2^n\cdot S_0=2^{n-1}\cdot (n-1)+2^{n-1}\cdot2=2^{n-1}\cdot (n+1).$ which is the same as above.  To simplify the process of finding the largest square $S_n$ we can note that if $n-1$ is odd then $n+1$ must be exactly divisible by an odd power of $2$.  However, this means $n+1$ is even but it cannot be.  Thus $n-1$ is even and $n+1$ is a large even square.  The largest even square $< 1000$ is $900$ so $n+1= 900 => n= \boxed{899}$

 





\textbf{Solution 3 (Finding patterns and using recursion)}

At first, this problem looks kind of daunting, but we can easily solve this problem by finding patterns, recursion and algebraic manipulations.

 We first simplify all the messy notation in the $S_n$ term. Note that the problem asks us to find the smallest value of $n<1000$ such that there exists an integer $k$ that satisfies

 
 \[k^2 = g(2) + g(4) + \cdots + g(2^n)\].

 Since there is no obvious way to approach this problem, we start by experimenting with small values of $n$ to evaluate some $S_n$.

 We play with these values:

 
 \[S_1 = g(2) = 2\]

 
 \[S_2 = g(2) + g(4) = 2+4 = 6\]

 
 \[S_3 = g(2) + g(4) + g(6) + g(8) = 16\]

 
 \[S_4 =  g(2) + g(4) + g(6) + g(8) + g(10) +g(12)+g(14)+g(16) = 40\]

 We are certainly not going to expand all of this out... so let's look for patterns from these $4$ values!

 Using a little bit of ingenuity, we note that

 
 \[S_2 = 2+4 = S_1 + 4\]

 
 \[S_3 = 2+4+2+8 = 8+8 = S_2 + S_1 + 8\]

 
 \[S_4 = 2+4+2+8+2+4+2+16 = S_3 + S_2 + S_1 + 16\]

 Aha! We see powers of two in each of our terms! Therefore, we can say that 

 
 \[S_2 = S_1 + 2^2\]

 
 \[S_3 = S_2+S_1 + 2^3\]

 
 \[S_4 = S_3 + S_2 + S_1 + 2^4\]

 Hooray! We have a recursion! Realistically, we would want to prove that the recursion works, but I currently don't know how to prove it. 

 On the actual AIME, go with whatever patterns you see, because most likely those are the patterns that the test-makers want the students to see.

 So we may generalize a formula for $S_n$:

 
 \[S_n = 2^n + S_{n-1} + S_{n-2} + \cdots + S_2 + S_1\]

 Uh oh... this formula is not in closed form. Looks like we'll have to use our recursion to develop one manually. We do so by using our recursion for $S_{n-1}$:

 
 \[S_n = 2^n + (2^{n-1} + S_{n-2} + S_{n-3} + \cdots + S_2 + S_1) + S_{n-2} + S_{n-3} + \cdots + S_2 + S_1\]

 
 \[S_n = 2^n + 2^{n-1} + 2 (S_{n-2} + S_{n-3} + \cdots + S_2 + S_1\]

 Let's simplify a bit further, where we use our recursion for $S_{n-2}$.

 
 \[S_n = 2^n + 2^{n-1} +(2S_{n-2}) + 2(S_{n-3} + S_{n-4} + \cdots + S_2 + S_1)\]

 
 \[S_n = 2^n + 2^{n-1} + 2(2^{n-2}) + 2(S_{n-3} + S_{n-4} + \cdots + S_2 + S_1) + 2(S_{n-3} + S_{n-4} + \cdots + S_2 + S_1)\]

 
 \[S_n = 2^n + 2^{n-1} + 2^{n-1} + 4(S_{n-3} + S_{n-4} + \cdots + S_2 + S_1)\]

 We now see a pattern! Using the exact same logic, we can condense this whole messy thing into a closed form:

 
 \[S_n = 2^n + \underbrace{2^{n-1}}_{n-2} + 2^{n-2}(S_1)\]

 
 \[S_n = 2^n + 2^{n-1}(n-2) + 2^{n-1}\]

 
 \[S_n = 2^n + 2^{n-1}(n-1)\]

 
 \[S_n = 2^{n-1}(2 + (n-1)\]

 
 \[S_n = 2^{n-1}(n+1)\]

 We have our closed form, so now we can find the largest value of $n$ such that $S_n$ is a perfect square.

 In order for $S_n$ to be a perfect square, we must have $n-1$ even and $n+1$ be a perfect square.

 Since $n<1000$, we have $n+1 < 1001$. We first try $n+1 = 31^2 = 961$(since it is the smallest square below $1000$, which gives us $n=960$. But $n-1$ isn't even, so we discard this value.

 Next, we try the second smallest value, which is $n = 30^2 = 900$, which tells us that $n=899$. $n-1$ is indeed even, and $n+1$ is a perfect square, so the largest value of $n$ such that $S_n$ is a perfect square is $899$.

 Our answer is $\boxed{899}$.

 -FIREDRAGONMATH16

 





\textbf{Solution 4}

First, we set intervals. Say that a number $N$ falls strictly within $[2^k, 2^{k+1}]$. 

 $N<2^{k+1}=2^k+2^k$

 We can generalize this:

 If a number is in the form $N=2^k+2^{R}O$ where $O$ is a positive odd number, $R<k$:

 $N<2^{k+1}=2^k+2^k\Longrightarrow O<2^{k-R}$ so there are $2^{k-R-1}$ numbers that satisfy this form. For all numbers that satisfy this form, $g(N)=2^R$. The sum of $g(N)$ for all $N$ in this form and interval is $2^{k-1}$. $R$ can vary between $1$ and $k-1$, so the total sum is $\underbrace{2^{k-1}+2^{k-1}\cdots 2^{k-1}}_{k-1}=(k-1)2^{k-1}$. The domain of the function we are trying to find is all even integers in the interval  $[2^1, 2^n]$ so there are $n-1$ values of $k$. Now we have the sum $\sum_{k=1}^{n-1}{(k-1)2^{k-1}}=\sum_{i=1}^{n-2}{k\cdot2^{k}}$. However, we did not consider powers of two yet(since our interval was strictly between powers of 2), so we have to add $\sum_{k=1}^{n}{2^k}$. Note that $\sum_{i=1}^{n-2}{k\cdot2^{k}}=(2^1+2^2\cdots 2^{n-2})+(2^2+2^3\cdots 2^{n-2})\cdots +(2^{n-3}+2^{n-2})+2^{n-2}=2^1(2^{n-2}-1)+2^2(2^{n-3}-1)\cdots +2^{n-3}(2^2-1)+2^{n-2}(2^1-1)=(n-2)(2^{n-1})-\sum_{k=1}^{n-2}{2^k}=(n-2)(2^{n-1})-2(2^{n-2}-1)=(n-3)(2^{n-1})+2$. Adding $\sum_{k=1}^{n}{2^k}=2(2^n-1)$, we get $(n+1)(2^{n-1})$. If this sum is a perfect square, $n\not\equiv0\pmod2$, since that would make $2^{n-1}$ not a perfect square, and $n+1$ odd so it cannot contribute a factor of 2 to make the power of 2 a perfect square.

 We want the least odd number less than $1000$ such that $(n+1)$ is an even perfect square. The greatest even square less than $1000$ is $30^2=900$ so $n+1=900 \Longrightarrow n=\boxed{899}$

 




	\item (\textit{AIME I 2006}) A tripod has three legs each of length 5 feet. When the tripod is set up, the angle between any pair of legs is equal to the angle between any other pair, and the top of the tripod is 4 feet from the ground. In setting up the tripod, the lower 1 foot of one leg breaks off. Let $h$ be the height in feet of the top of the tripod from the ground when the broken tripod is set up. Then $h$ can be written in the form $\frac m{\sqrt{n}},$ where $m$ and $n$ are positive integers and $n$ is not divisible by the square of any prime. Find $\lfloor m+\sqrt{n}\rfloor.$ (The notation $\lfloor x\rfloor$ denotes the greatest integer that is less than or equal to $x.$)


 
 \textbf{Answer: }183



\textbf{Solution 1}


                    \begin{figure}[H]
                        \centering
                        \includegraphics[width=0.5\textwidth]{Solution/AIME_2006_I_14_Solution_Figure_1.png}
                    \end{figure}
                    We will use $[...]$ to denote volume (four letters), area (three letters) or length (two letters).

 Let $T$ be the top of the tripod, $A,B,C$ are end points of three legs. Let $S$ be the point on $TA$ such that $[TS] = 4$ and $[SA] = 1$. Let $O$ be the center of the base \href{https://artofproblemsolving.com/wiki/index.php?title=Equilateral\_triangle}{equilateral triangle} $ABC$. Let $M$ be the \href{https://artofproblemsolving.com/wiki/index.php?title=Midpoint}{midpoint} of segment $BC$. Let $h$ be the distance from $T$ to the triangle $SBC$ ($h$ is what we want to find).

 We have the volume ratio $\frac {[TSBC]}{[TABC]} = \frac {[TS]}{[TA]} = \frac {4}{5}$.

 So $\frac {h\cdot [SBC]}{[TO]\cdot [ABC]} = \frac {4}{5}$.

 We also have the \href{https://artofproblemsolving.com/wiki/index.php?title=Area\_ratios}{area ratio} $\frac {[SBC]}{[ABC]} = \frac {[SM]}{[AM]}$.

 The triangle $TOA$ is a $3-4-5$ \href{https://artofproblemsolving.com/wiki/index.php?title=Right\_triangle}{right triangle} so $[AM] = \frac {3}{2}\cdot[AO] = \frac {9}{2}$ and $\cos{\angle{TAO}} = \frac {3}{5}$.

 Applying \href{https://artofproblemsolving.com/wiki/index.php?title=Law\_of\_Cosines}{Law of Cosines} to the triangle $SAM$ with $[SA] = 1$, $[AM] = \frac {9}{2}$ and $\cos{\angle{SAM}} = \frac {3}{5}$, we find:

 $[SM] = \frac {\sqrt {5\cdot317}}{10}.$Putting it all together, we find $h = \frac {144}{\sqrt {5\cdot317}}$.

 $\lfloor 144+\sqrt{5 \cdot 317}\rfloor =144+ \lfloor \sqrt{5 \cdot 317}\rfloor =144+\lfloor \sqrt{1585} \rfloor =144+39=\boxed{183}$.





\textbf{Solution 2}

We note that $AO=3$. From this we can derive that the side length of the equilateral is $3\sqrt{3}$. We now use 3D coordinate geometry. 
 \[A = (0,0,0)\] 
 \[B = (3\sqrt{3},0,0)\] 
 \[C = (\frac{3\sqrt{3}}{2}, \frac{9}{2}, 0)\] 
 \[T = (\frac{3\sqrt{3}}{2}, \frac{3}{2}, 4)\] 
 \[S= (\frac{3\sqrt{3}}{10}, \frac{3}{10}, \frac{4}{5})\]We know three points of plane $SCB$ hence we can write out the equation for the plane. Plane $SCB$ can be expressed as 
 \[4\sqrt{3}x+4y+39z-36=0.\]

 Applying the distance between a point and a plane formula. 

 
 \[\frac{ax+by+cz+d}{\sqrt{a^{2}+b^{2}+c^{2}}} = \frac{4\sqrt{3} \cdot \frac{3\sqrt{3}}{2} + 4\cdot \frac{3}{2} + 39 \cdot 4 -36}{\sqrt{(4\sqrt{3})^2+4^2+39^2}} = \frac{144}{\sqrt{1585}}\]

 
 \[\lfloor m+\sqrt{n}\rfloor = \lfloor 144+\sqrt{1585}\rfloor = 183\]

 Solution by SimonSun

 


 





\textbf{Solution 3 (Cosine Law & Pythagorean Bash)}

Diagram borrowed from Solution 1

 
                    \begin{figure}[H]
                        \centering
                        \includegraphics[width=0.5\textwidth]{Solution/AIME_2006_I_14_Solution_Figure_1.png}
                    \end{figure}
                    Apply \href{https://artofproblemsolving.com/wiki/index.php?title=Pythagorean\_Theorem}{Pythagorean Theorem} on $\bigtriangleup TOB$ yields
 \[BO=\sqrt{TB^2-TO^2}=3\]Since $\bigtriangleup ABC$ is equilateral, we have $\angle MOB=60^{\circ}$ and 
 \[BC=2BM=2(OB\sin MOB)=3\sqrt{3}\]Apply Pythagorean Theorem on $\bigtriangleup TMB$ yields
 \[TM=\sqrt{TB^2-BM^2}=\sqrt{5^2-(\frac{3\sqrt{3}}{2})^2}=\frac{\sqrt{73}}{2}\]Apply \href{https://artofproblemsolving.com/wiki/index.php?title=Law\_of\_Cosines}{Law of Cosines} on $\bigtriangleup TBC$ we have
 \[BC^2=TB^2+TC^2-2(TB)(TC)\cos BTC\]
 \[(3\sqrt{3})^2=5^2+5^2-2(5)^2\cos BTC\]
 \[\cos BTC=\frac{23}{50}\]Apply Law of Cosines on $\bigtriangleup STB$ using the fact that $\angle STB=\angle BTC$ we have
 \[SB^2=ST^2+BT^2-2(ST)(BT)\cos STB\]
 \[SB=\sqrt{4^2+5^2-2(4)(5)\cos BTC}=\frac{\sqrt{565}}{5}\]Apply Pythagorean Theorem on $\bigtriangleup BSM$ yields
 \[SM=\sqrt{SB^2-BM^2}=\frac{\sqrt{1585}}{10}\]Let the perpendicular from $T$ hits $SBC$ at $P$. Let $SP=x$ and $PM=\frac{\sqrt{1585}}{10}-x$. Apply Pythagorean Theorem on $TSP$ and $TMP$ we have
 \[TP^2=TS^2-SP^2=TM^2-PM^2\]
 \[4^2-x^2=(\frac{\sqrt{73}}{2})^2-(\frac{\sqrt{1585}}{10}-x)^2\]Cancelling out the $x^2$ term and solving gets $x=\frac{181}{2\sqrt{1585}}$.

 Finally, by Pythagorean Theorem,
 \[TP=\sqrt{TS^2-SP^2}=\frac{144}{\sqrt{1585}}\]so $\lfloor m+\sqrt{n}\rfloor=\boxed{183}$

 




	\item (\textit{AIME I 2006}) Given that a sequence satisfies $x_0=0$ and $|x_k|=|x_{k-1}+3|$ for all integers $k\ge 1,$ find the minimum possible value of $|x_1+x_2+\cdots+x_{2006}|.$


 
 \textbf{Answer: }027



\textbf{Solution 1}

Suppose $b_{i} = \frac {x_{i}}3$.We have
 \[\sum_{i = 1}^{2006}b_{i}^{2} = \sum_{i = 0}^{2005}(b_{i} + 1)^{2} = \sum_{i = 0}^{2005}(b_{i}^{2} + 2b_{i} + 1)\]So
 \[\sum_{i = 0}^{2005}b_{i} = \frac {b_{2006}^{2} - 2006}2\]Now
 \[\sum_{i = 1}^{2006}b_{i} = \frac {b_{2006}^{2} + 2b_{2006} - 2006}2\]Therefore
 \[\left|\sum_{i = 1}^{2006}b_{i}\right| = \left|\frac {(b_{2006} + 1)^{2} - 2007}2\right|\geq \left|\frac{2025 - 2007}{2}\right| = 9\]This lower bound can be achieved if we use $b_1 = -1$, $b_2 = 0$, $b_3 = -1$, $b_4 = 0$, and so on until $b_{1962} = 0$, after which we let $b_k = b_{k - 1} + 1$ so that $b_{2006} = 44$. So
 \[\left|\sum_{i = 1}^{2006}x_{i}\right|\geq \boxed{027}\]

 





\textbf{Solution 2}

First, we state that \href{https://artofproblemsolving.com/wiki/index.php?title=Iff}{iff} $x_{i - 1}\ge0$, $|x_i| = |x_{i - 1}| + 3$ and iff $x_{i - 1} < 0$, $|x_i| = |x_{i - 1}| - 3$. Now suppose $x_i = x_j$ for some $0\le i < j\le2006$. Now, this means that $|x_i| = |x_j|$, and so the number of positive numbers in the set $\{x_i,x_{i + 1},\ldots,x_{j - 1}\}$ equals the number of negative numbers. Now pair the numbers in this list up in the following way: Whenever a positive and a negative number are adjacent in this progression, pair them up and remove them from this list. We claim that every pair will sum to -3.

 If the positive number comes first, then the negative number will have a magnitude three greater, so this is true. If the negative number comes first, then the positive number will have magnitude three smaller, and this will also be true. Now let us examine what happens when we remove those two from the sequence. WLOG, let the numbers be $x_k$ and $x_{k + 1}$. Since one is positive and the other is negative, $|x_{k + 2}| = |x_{k + 1}|\pm3 = |x_k|\pm3\mp3 = |x_k| = |x_{k - 1} + 3|$. So the new sequence works under the same criteria as the old one. In this way, we can pair all of the numbers up in this subsequence so the sums of the pairs are -3. Thus, the average of these numbers will be -3/2 for all subsequences that start and end with the same number (not including one of those).

 Now, take all of the repeating subsequences out of the original sequence. The only thing that will be left will be a sequence $\{0,3,6,9,\cdots,3k\}$ for some even $k$. Since we started with 2006 terms, we removed $2006 - k$ (an even number) with an average of -3/2. Thus, the sum of both this remaining sequence and the removed stuff is $(2006 - k)( - 3/2) + \sum_{i = 1}^k3k = (3/2)(k - 2006 + k(k + 1)) = 3/2(k^2 + 2k - 2006)$. This must be minimized, so we find the roots: $k^2 + 2k = 2006\implies (k + 1)^2 = 2007$ and $44^2 = 1936 < 2007 < 2025 = 45^2$. Plugging in $k = 44$ yields $(3/2)(2025 - 2007) = 27$ (and $k = 42$ yields $- 237$, a worse result). Thus, $\fbox{027}$ is the closest to zero this sum can get.

 





\textbf{Solution 3}

We know $|x_k| = |x_{k - 1} + 3|$.  We get rid of the absolute value by squaring both sides: ${x_k}^2 = {x_{k - 1}}^2 + 6{x_{k - 1}} + 9\Rightarrow {x_k}^2 - {x_{k - 1}}^2 = 6{x_{k - 1}} + 9$. So we set this up:

 
 \begin{align*} {x_1}^2 - {x_0}^2 & = 6{x_0} + 9 \\ {x_2}^2 - {x_1}^2 & = 6{x_1} + 9 \\ & \vdots \\ {x_{2007}}^2 - {x_{2006}}^2 & = 6{x_{2006}} + 9 \end{align*}
There are $2007$ equations. Sum them. We get:${x_{2007}}^2 = 6\left(x_1 + x_2 + \cdots + x_{2006}\right) + 9\cdot{2007}$

 So $|x_1 + x_2 + \cdots + x_{2006}| = \frac 16 \left|{x_{2007}}^2 - 9\cdot{2007}\right|$

 We know $3\ |\ x_{2007}$ and we want to minimize $\left|{x_{2007}}^2 - 9\cdot{2007}\right|$, so $x_{2007}$ must be $3\cdot{45}$ for it to be minimal ($45^2 = 2025$ which is closest to $2007$). We can achieve this with $x_k=3k$ till $x_{45}=135$ and then alternating $x_{46}=-138$, $x_{47}=135$ and so on ... Then $x_{2k}=-138$ and $x_{2k+1}=135$ for all $k>22$. Since $2007$ is odd, we have $x_{2007}=135$.

 This means that $|x_1 + x_2 + \cdots + x_{2006}| = \frac 16 \left|9(2025 - 2007)\right| = \boxed{027}$

 





\textbf{Solution 4}

Playing around with a couple numbers, we see that we can generate the sequence $0, 3, -6, 3, -6, \cdots$, and we can also generate the sequence $3, 6, 9, 12, \cdots$ after each $-6$ value. Thus, we will apply this to try and find some bounds. We can test if the first $1000$ pairs of numbers each sum up to $-3$, and the rest form an arithmetic sequence, if the first $990$ pairs sum up to $-3$, and so on. When we get to $980$, we find that $980(-3) + 3 + 6 + \cdots + 3\cdot 46 = 303$. If we shift the number of pairs up by $1$, we get $981(-3) + 3 + 6 + \cdots + 3\cdot 44 = \boxed{027}$.

 





\textbf{Solution 5}

We will work our way from $|x_0|$ to $|x_0+x_1|$ to $|x_0+x_1+x_2+\dots+x_{2006}|$. Let the sum $|x_0+x_1+x_2+\dots+x_i|=S_i$. 

 Seeing as the value for the sum we want is always nonnegative, the best pseudo-strategy would be to stay as close to $0$ as possible as we increase $i$ to eventually get to $i=2006$. It turns out that a greedy algorithm works here. Let us start with some smaller values of $i$. 

 Note that we can describe $x_{k+1}$ in terms of $x_k$ in the following way: take $x_k$, add $3$, and either multiply by $-1$ or not. 

 We know that $S_0=0$. When we get to $S_1$, we add in $x_1$, which is either $3$ or $-3$. It does not make a difference which one we choose, so we can choose $3$ for convenience. Now, $S_1=3$. $x_2$ is either $\pm6$; adding $6$ would take us farther away from $0$, but adding $-6$ is an okay move. Thus, let $x_2=-6$, so $S_2=-3$. 

 $x_3$ is $\pm3$. Adding $3$ is clearly the right choice here, which takes us to $0$. Thus, $x_3=3$ and $S_3=0$. 

 Let us look at how this greedy algorithm performs in general. 

 $\begin{tabular}[t]{c|ccccccc}S_i&3a&-3&3a-3&-6&3a-6&-9&\cdots\\x_i&3a&-3a-3&3a&-3a-3&3a&-3a-3&\cdots\\ \end{tabular}$ (Error compiling LaTeX. Unknown error\_msg)

 (the table doesn't work; if you desire, please to go \href{https://artofproblemsolving.com/texer/zzyacvnp}{https://artofproblemsolving.com/texer/zzyacvnp} to view)

 The sums alternate between separate arithmetic sequences. One of them is of particular interest to us: specifically, the one that goes $3a,3a-3,3a-6,\cdots$, because it will eventually become $0$. One can then derive using simple algebra how long it will take to reach zero, depending on $a$; it turns out that it takes $2a+1$ for each cycle to complete. We can then see that $S_{k^2-1}=0$ for positive integers $k$; thus, $S_{1935}=0$. We can then go through our algorithm, and it turns out that $S_{2006}=\boxed{027}$. 

 




\end{enumerate}